\subsection{Database design}

The database was designed for the PostgreSQL DBMS selected in Activity A1.1~\ref{act:1.1} and for the Django-based backend defined in the technology stack. The schema supports the complete workflow of the application: project creation, query definition, YouTube data collection, API request traceability, data processing and visualization. In this design, Django model classes define the application entities and the ORM maps them to PostgreSQL tables, so the database structure remains close to the backend implementation.

The design also responds to Activity A1.3~\ref{act:1.3}, which requires storage for videos, comments, replies, user projects, API queries and processing metadata. For methodological transparency, the schema does not only store the final processed data. It also stores collection runs, request parameters, quota consumption, pagination information, raw API payloads and processing outputs, because these records describe the conditions under which the dataset was obtained and transformed.

The complete table definitions and ER diagram are in \href{https://github.com/LONG-TERM-EFFECTS-OF-SUFFERING/thesis/tree/main/docs/database_schema.md}{\texttt{database\_schema.md}}. The following sections summarize the design criteria, implementation model, main entities, relationships and integrity controls of the database schema.

\subsubsection{Design criteria}

The schema was defined using the following criteria:

\begin{itemize}
    \item Preserve \textbf{traceability} of each collection session through explicit records of the effective query, execution time, API endpoints, request parameters, quota consumption and pagination tokens.

    \item Store YouTube resources in a \textbf{normalized relational structure} that separates channels, videos, comments, request logs, processing outputs and project management data.

    \item Represent \textbf{replies} in the same table as top-level comments using a self-reference, instead of creating a duplicated table structure for replies.

    \item Preserve \textbf{reproducibility} by keeping original YouTube identifiers, raw JSON payloads, timestamps, query parameters and processing metadata.

    \item Support later \textbf{analysis and visualization} through sentiment results, derived metrics and stable relationships between projects, collection runs, videos and comments.
\end{itemize}

\subsubsection{Implementation model}

The current model uses Django's built-in authentication user through \texttt{settings.\allowbreak AUTH\_\allowbreak USER\_\allowbreak MODEL}. With Django's default configuration, this means that user accounts are stored in \texttt{auth\_user}, not in a custom table. Application-specific entities use UUID primary keys generated by Django, while foreign keys that point to the user model follow the primary key type of the configured Django user.

PostgreSQL \texttt{JSONB} columns are represented in Django with \texttt{JSONField}. These fields are used when the application must store flexible data whose structure may vary between requests, such as structured query parameters, raw API responses, request payloads and summary outputs. Timestamp columns such as \texttt{created\_at} and \texttt{updated\_at} are maintained by Django using automatic timestamp fields. Enumerated values such as statuses, query sources, comment levels and sentiment labels are represented with Django choices.

\begin{table}[H]
    \centering
    \small
    \caption{Database tables organized by functional area}
    \label{tab:database_tables_by_area}
    \begin{tabularx}
        {\textwidth}{>{\raggedright\arraybackslash}p{3.2cm}>{\raggedright\arraybackslash}p{5.0cm}X}
        \toprule
        \textbf{Functional area}    & \textbf{Tables}                                                                                                                                                                                                                                                       & \textbf{Purpose}                                                                                                                                                                                               \\
        \midrule
        User and project management & \texttt{auth\_user}, \texttt{research\_projects}, \texttt{saved\_queries}                                                                                                                                                                                             & Stores users, research workspaces and reusable query definitions. Saved queries keep both the natural language request and the structured YouTube API parameters generated manually or through the LLM module. \\
        \midrule
        Collection traceability     & \texttt{collection\_runs}, \texttt{api\_request\_logs}                                                                                                                                                                                                                & Documents each collection execution, including status, counts, timestamps, quota usage, endpoint names, request parameters, pagination tokens and API response metadata.                                       \\
        \midrule
        YouTube resources           & \texttt{youtube\_channels}, \texttt{youtube\_videos}, \texttt{collection\_\allowbreak run\_\allowbreak videos}, \texttt{video\_\allowbreak statistics\_\allowbreak snapshots}, \texttt{youtube\_comments}, \texttt{collection\_\allowbreak run\_\allowbreak comments} & Stores normalized platform data and the linking records that explain which videos and comments were discovered or collected during each run.                                                                   \\
        \midrule
        Processing and analysis     & \texttt{processing\_runs}, \texttt{comment\_sentiments}, \texttt{derived\_metrics}                                                                                                                                                                                    & Stores executions of the processing pipeline, model-based sentiment results and aggregated values used by charts and summary views.                                                                            \\
        \bottomrule
    \end{tabularx}
\end{table}

\subsubsection{Main entities}

\paragraph{User and project management}

The project uses Django's \texttt{auth\_user} table for account management. The table \texttt{research\_projects} stores the research workspaces created by users, and each project belongs to one owner. The pair \texttt{owner\_user\_id} and \texttt{name} is unique so that a user cannot create two projects with the same name. The table \texttt{saved\_queries} stores reusable query definitions for each project. It includes the original natural language prompt (if there was one), the query source, date filters, regional and language filters, structured query parameters, LLM model metadata (if the query was generated by a LLM) and optional notes. This design connects the user-facing query process with the exact parameters later sent to the YouTube Data API.

\paragraph{Collection traceability}

The table \texttt{collection\_runs} represents one execution of the data collection workflow. It stores the project, the saved query used when applicable, the user who initiated the run, execution status, requested comment limit, start and finish timestamps, API request counts, quota units consumed, collected item counts and error information. The table \texttt{api\_request\_logs} stores the detailed request history for each run. Each log records the request sequence, endpoint name, request parameters, \texttt{part} and \texttt{fields} parameters, pagination tokens, HTTP status, quota cost, number of returned items and possible errors. These tables make the collection process auditable instead of treating it as a black box.

\paragraph{YouTube resources}

The tables \texttt{youtube\_channels} and \texttt{youtube\_videos} store normalized metadata for channels and videos. Each channel keeps its original \texttt{youtube\_channel\_id}, title, description, country, publication date, statistics and raw payload. Each video keeps its original \texttt{youtube\_video\_id}, channel reference, title, description, publication date, duration, language metadata, category, tags, live broadcast status and raw payload. Original YouTube identifiers are unique business keys, while UUID values are used internally for application relationships.

Because the same video can appear in multiple collection runs, the many-to-many relationship between runs and videos is stored in \texttt{collection\_run\_videos}. This link table also stores discovery context, such as page number and discovery timestamp. Video statistics are stored in \texttt{video\_statistics\_snapshots} instead of only inside \texttt{youtube\_videos}, because values such as views, likes and comment counts can change over time.

The table \texttt{youtube\_comments} stores both top-level comments and replies. Each row belongs to one video and keeps the original \texttt{youtube\_comment\_id}, author channel identifier, displayed text, original text, like count, reply count, publication timestamps, collection timestamps and raw payload. Replies are represented by \texttt{parent\_comment\_id}, a self-referential foreign key that points from a reply to its parent top-level comment. The \texttt{comment\_level} field distinguishes \texttt{top\_level} comments from \texttt{reply} comments, and a check constraint should enforce that replies have a parent while top-level comments do not. The table \texttt{collection\_run\_comments} links comments to the run that collected them and stores page and thread-position context.

\paragraph{Processing and analysis}

The table \texttt{processing\_runs} stores each execution of the processing pipeline. It links processing output to a project and the collection run used as input. It also records status, rows read, rows written, summary payload, error messages and execution timestamps. The table \texttt{comment\_sentiments} stores the sentiment result produced for each comment within a processing run, including model name, label, numeric score, confidence, language code and raw model output. The table \texttt{derived\_metrics} stores aggregated analytical values by scope, such as project-level, collection-run-level, channel-level, video-level and comment-level. These records provide the structured inputs used by the visualization layer.

\subsubsection{Relationship summary}

The main entity relationships are the following:

\begin{itemize}
    \item One Django user can own many research projects.

    \item One research project can contain many saved queries, collection runs and processing runs.

    \item One saved query can be used by many collection runs and a collection run need to be created with a saved query.

    \item One collection run can generate many API request logs.

    \item One collection run can be linked to many videos through \texttt{collection\_run\_videos}, and one video can appear in many collection runs.

    \item One channel can publish many videos.

    \item One video can have many statistics snapshots and many comments.

    \item One top-level comment can have many replies through the self-reference \texttt{parent\_\allowbreak comment\_\allowbreak id}.

    \item One collection run can be linked to many comments through \texttt{collection\_\allowbreak run\_\allowbreak comments}, and one comment can appear in many collection runs.

    \item One processing run can generate many sentiment records and many derived metrics.

    \item One comment can have many sentiment records when it is analyzed in different processing runs or by different models.
\end{itemize}

\subsubsection{Integrity and reproducibility controls}

Several integrity controls are included to keep the stored data consistent. Unique constraints prevent duplicate project names per user, duplicate query names per project, duplicate external YouTube identifiers, duplicate video links within the same collection run and duplicate comment links within the same collection run. The request log sequence is unique inside each collection run, so the request order can be reconstructed. Date checks should ensure that finish timestamps are not earlier than start timestamps and that saved query date ranges are valid.

The reproducibility controls are equally important. Collection tables store the exact request parameters, quota costs and pagination tokens used to obtain the data. YouTube resource tables preserve the original platform identifiers and raw payloads. Processing tables store model names, summary payloads and raw model output. Together, these records allow a later reviewer to understand not only what data were collected, but also how the data moved from API requests to stored records and analytical results.
